\scanpage{8}

% 1, continuation
avec les $\sigma_i$ opérant ss pt fixe.

% 2–5, proposition
\uline{Prop} le foncteur de la catégorie isotopique des \inserted{variétés de dim $n$} $X$ filtrés comme dessus, vers les $G_n$ ensembles finis (où les $\sigma_i$ opèrent ss pt fixe) est \uline{pl. fidèle}.

% 6–11, essential image
Déterminer l'image essentielle (\unsure{admettre|mettre|moyennant} l'hyp. de Poincaré au besoin) et montrer qu'elle est plus grande que celle correspondant aux complexes cellulaires se réalisant suivant une variété topologique\ldots.

% 12–16, combinatorial variety
\uline{$n$-Variété combinatoire} \inserted{\uline{finie}}, définie par un ens \inserted{fini} $D_n$ muni d'autom. invol. sans pts fixes $\sigma_i$ $(0\leq i\leq n)$ tels que deux $\sigma_i$ \unsure{:} non consécutifs commutent --- i.e. $D_n$ est un ens. à op. sous $G_n$, avec les $\sigma_i$ $(0\leq i\leq n)$ op. sans pts fixes.

% 17–20, combinatorial complex
\uline{$n$-complexe combinatoire} \uline{fini} Donné par un système $D_n,D_{n-1},\ldots,D_1,\underset{\substack{\downarrow\\\text{``spécial''}}}{D_0}$ d'ens \inserted{finis}, chaque \inserted{$D_i\ (1\leq i\leq n)$} étant un \deleted{\unsure{$D_i$}} $G_{i-1}$-ensemble (\deleted{\unsure{à él\textsuperscript{ts}|i.e.}} les gén. $\sigma_\alpha$ \unsure{op.} sans pts fixes) et des applications

% 21, sequence of applications
\[
 D_n\to D_{n-1}\longrightarrow D_{n-2}\longrightarrow\ldots\longrightarrow D_1\longrightarrow D_0
\]

% 22–23, compatibility, continues
telles que $\forall\,2\leq i\leq n$, $D_i\to D_{i-1}$ soit un $G_{i-2}$-morphisme, et que pour $\deleted{\unsure{2\leq j\geq{}|0\leq j\leq{}}}i\geq j+2\geq j\geq0$,

% lower left
% archive sheet number: 7
